\chapter{Weitere Themen und Methoden: Ideen und Ergänzungen}\label{ch:Appendix_Weitere_Themen}

\begin{myremark}
Dieser Anhang listet inhaltliche Themen und methodische Ansätze, die das
KI-Skript sinnvoll ergänzen könnten. Es handelt sich um Vorschläge und
Ideen für eine Weiterentwicklung des Moduls.
\end{myremark}

\section*{A. Inhaltliche Themen}

\subsection*{A.1 Technische Vertiefungen}

\begin{itemize}
    \item \textbf{Transformerarchitektur} (vereinfacht): Wie funktioniert
          der Attention-Mechanismus? Warum können Transformer so gut auf
          langen Texten arbeiten? Mögliche Umsetzung: interaktive Visualisierung
          mit \url{https://bbycroft.net/llm}.
    \item \textbf{Embeddings und Vektorraummodelle}: Wie repräsentiert ein LLM
          Wörter als Vektoren? Übung: Ähnliche Wörter im Vektorraum finden.
    \item \textbf{Retrieval Augmented Generation (RAG)}: Wie können LLMs auf
          aktuelle oder private Dokumente zugreifen? Praktische Demo mit einem
          einfachen RAG-System (z.\,B. LlamaIndex).
    \item \textbf{Multimodalität}: Wie verarbeiten moderne Modelle gleichzeitig
          Text, Bild, Audio und Video? Übung mit Gemini oder GPT-4o.
    \item \textbf{Quantisierung und Edge-AI}: Wie können Modelle auf kleineren
          Geräten (Smartphone, Laptop) ohne Cloud ausgeführt werden?
          Demo: Ollama lokal auf Schullaptop.

\end{itemize}

\subsection*{A.2 Gesellschaftliche Vertiefungen}

\begin{itemize}
    \item \textbf{KI und Demokratie}: Auswirkungen von KI-generierten
          Desinformationskampagnen auf Wahlen. Fallstudie zu einem konkreten
          Ereignis (z.\,B. Deepfakes im US-Wahlkampf).
    \item \textbf{KI-Regulierung}: Vergleich EU AI Act vs. Schweizer Ansatz
          vs. US-amerikanischer Ansatz. Planspiel: Sie sind Parlamentsmitglied
          und müssen ein KI-Gesetz verabschieden.
    \item \textbf{Überwachungskapitalismus}: Wie finanzieren sich kostenlose
          KI-Dienste? Diskussion des Geschäftsmodells (Daten = Währung).
    \item \textbf{Militärische KI}: Drohnen, autonome Waffen, LAWS (Lethal
          Autonomous Weapon Systems). Ethische Debatte.
    \item \textbf{KI in der Medizin}: Diagnose, Medikamentenentwicklung,
          personalisierte Medizin. Chancen und Risiken.
    \item \textbf{KI und Kreativität}: Kann eine KI wirklich kreativ sein?
          Philosophische Diskussion + Vergleich KI-generierte vs. menschliche Kunst.
    \item \textbf{Digitale Kluft}: Wer hat Zugang zu leistungsstarken KI-Tools?
          Reproduziert KI bestehende globale Ungleichheiten?
\end{itemize}

\subsection*{A.3 Lokale und aktuelle Bezüge (Schweiz)}

\begin{itemize}
    \item \textbf{Schweizer KI-Ökosystem}: ETH Zürich, EPFL, AI-Startups
          (z.\,B. Aleph Alpha Schweiz-Niederlassung). Besuch / Zoom-Call
          mit einer Fachperson.
    \item \textbf{Schweizer Volksabstimmungen und KI}: Könnten in Zukunft
          KI-Systeme den Abstimmungskampf beeinflussen?
          Simulationsübung.
    \item \textbf{KI und Schweizer Berufsbildung}: Welche Ausbildungsberufe
          verändern sich besonders? Befragung von Berufsbildnerinnen
          und Berufsbildnern.
\end{itemize}

\section*{B. Methodische Ergänzungen}

\subsection*{B.1 Interaktive Methoden}

\begin{itemize}
    \item \textbf{Blind-Test}: SuS lesen Texte (ein von Mensch geschriebener,
          einer von KI) und versuchen zu erkennen, welcher welcher ist.
          Anschliessend: Reflexion über die Qualität der Unterscheidung.
    \item \textbf{Turing-Test simulieren}: Eine Person chattet mit einer
          anderen Person und mit einem LLM gleichzeitig, ohne zu wissen,
          welches welches ist. Kann sie den Unterschied erkennen?
    \item \textbf{Jigsaw-Methode}: Jede Expertengruppe erarbeitet einen
          Aspekt (z.\,B. Technologie, Ethik, Wirtschaft, Recht) und lehrt
          dann die anderen Gruppen.
    \item \textbf{Zukunftswerkstatt}: SuS entwickeln in Gruppen Szenarien
          für das Jahr 2035 (Utopie / Dystopie) und präsentieren diese.
    \item \textbf{Fishbowl-Diskussion}: Innerer Kreis debattiert, äusserer
          beobachtet und gibt Feedback. Geeignet für ethische Themen.
    \item \textbf{World Café}: Verschiedene Tische mit je einem Aspekt
          von KI, SuS rotieren und ergänzen die Erkenntnisse der Vorgruppe.
\end{itemize}

\subsection*{B.2 Projekte und Langzeitaufgaben}

\begin{itemize}
    \item \textbf{KI-Tagebuch}: SuS führen während des Moduls ein Tagebuch,
          in dem sie jeden KI-Kontakt (bewusst oder unbewusst) dokumentieren.
          Abschlusspräsentation: Wie viel KI ist in meinem Alltag?
    \item \textbf{Mini-Maturaarbeit KI}: SuS untersuchen ein konkretes
          KI-Anwendungsfeld vertieft und präsentieren Ergebnisse.
    \item \textbf{KI-Produkt entwickeln}: Gruppen entwickeln ein Konzept
          für ein KI-basiertes Produkt, das ein echtes Problem löst.
          Elevator-Pitch vor der Klasse.
    \item \textbf{KI-Zeitung}: Die Klasse erstellt gemeinsam eine
          Schülerzeitung komplett mit Unterstützung von KI-Tools
          und reflektiert anschliessend den Prozess.
\end{itemize}

\subsection*{B.3 Externe Ressourcen und Gastlektionen}

\begin{itemize}
    \item \textbf{AI Literacy-Kurse}: Elements of AI (\url{https://www.elementsofai.com/de})
          --- kostenloser Online-Kurs, auf Deutsch verfügbar, für Sek-II-Niveau geeignet.
    \item \textbf{Gastlektionen}: Fachpersonen aus KI-nahen Berufen einladen
          (Startup, ETH, Versicherung, Journalismus).
    \item \textbf{Museum / Ausstellung}: \enquote{Exploring AI} oder vergleichbare
          Wanderausstellungen in der Schweiz.
    \item \textbf{Medienbegleitung}: NZZ, SRF Digital, The Decoder als laufende
          Informationsquellen für aktuelle KI-Entwicklungen.
\end{itemize}

\section*{C. Lernziele (Entwurf)}

\begin{itemize}
    \item Ich kann erklären, wie ein LLM auf einer grundlegenden Ebene funktioniert
          (Token, Training, RLHF).
    \item Ich kann mindestens drei kritische Aspekte von LLMs benennen und anhand
          von Beispielen erläutern.
    \item Ich kann aktuelle KI-Tools (LLMs, Bild-, Video-, Audio-Generatoren)
          kategorisieren und ihre Stärken und Schwächen beschreiben.
    \item Ich kann Prompts nach den Prinzipien des Prompt Engineerings formulieren
          und iterativ verbessern.
    \item Ich kann erklären, was ein KI-Agent ist, und konkrete Anwendungsbeispiele
          nennen.
    \item Ich kann die Datenschutzrelevanz der Nutzung von KI-Tools einschätzen
          und datenschutzkonforme Alternativen für den Schulbereich benennen.
    \item Ich kann die Auswirkungen von KI auf den Arbeitsmarkt in einem selbst
          gewählten Berufsfeld analysieren und präsentieren.
    \item Ich kann meine Ergebnisse kritisch reflektieren und die Grenzen
          von KI-generierten Antworten einschätzen.
\end{itemize}

\section*{D. Weiterführende Links und Ressourcen}

\begin{itemize}
    \item \textbf{LM Arena (Modell-Vergleich)}: \url{https://lmarena.ai}
    \item \textbf{Elements of AI (Grundkurs)}: \url{https://www.elementsofai.com/de}
    \item \textbf{The Decoder (News)}: \url{https://the-decoder.de}
    \item \textbf{AI Literacy Projekt}: \url{https://ailiteracy.eu}
    \item \textbf{Our World in Data --- KI}: \url{https://ourworldindata.org/artificial-intelligence}
    \item \textbf{Stanford AI Index Report} (jährlich):
          \url{https://aiindex.stanford.edu}
    \item \textbf{EU AI Act}: \url{https://artificialintelligenceact.eu}
    \item \textbf{Schweizerischer Datenschutz-Blog}: \url{https://www.datenschutz.ch}
    \item \textbf{OpenAI Tokenizer}: \url{https://platform.openai.com/tokenizer}
    \item \textbf{3Blue1Brown --- Neural Networks} (englisch, mit Untertiteln):
          \url{https://www.youtube.com/watch?v=aircAruvnKk}
    \item \textbf{LLM Visualizer (bbycroft)}: \url{https://bbycroft.net/llm}
\end{itemize}
